% !TEX root = ../main.tex

\chapter{Phase 3 - Market Insights}

\section{Objectives}

Phase 2 focused on collecting and analysing empirical data from Python-related freelance job postings. While the previous phase answered the individual
research questions separately, the objective of Phase 3 is to transform those findings into a coherent set of market insights that can support decision making.

Rather than introducing new analyses, this phase synthesises the results of the previous chapter to identify recurring market patterns, determine the most
valuable technical skills, compare different Python career paths, and highlight the technologies that appear most relevant for entry-level freelancers.

The resulting market insights provide the foundation for Phase 4, where they will be used to define the requirements of a market-oriented
portfolio project roadmap.

\section{Synthesis of the Market Analysis}

The previous phase investigated the analysed freelance market from five complementary perspectives. Each research question focused on a different
aspect of the collected job postings, ranging from the overall distribution of Python-related job categories to technology requirements and job compensation.

Individually, these analyses provide valuable information about specific market characteristics. However, considered together, they reveal a coherent
picture of the technical skills, technologies and market trends that are most relevant for Python freelancers. Table~\ref{tab:market_synthesis} summarises
the principal findings of each research question together with their practical interpretation.

\begin{table}[H]
    \centering
    \small
    \caption{Summary of the principal findings obtained during the market analysis.}
    \label{tab:market_synthesis}

    \begin{tabularx}{\textwidth}{
        c
        p{3.6cm}
        X
    }
        \toprule
        \textbf{RQ} &
        \textbf{Main Finding} &
        \textbf{Implication} \\
        \midrule

        RQ1 &
        Data Analysis is the largest job category. &
        Entry-level opportunities are concentrated primarily in Data Analysis. \\

        RQ2 &
        Pandas and NumPy dominate package requirements. &
        These libraries represent essential Python skills across multiple domains. \\

        RQ3 &
        Python, SQL and Excel appear most frequently. &
        A competitive freelancer should master the core technology stack rather than Python alone. \\

        RQ4 &
        Python is commonly requested together with SQL and Excel. &
        Portfolio projects should demonstrate the integration of complementary technologies. \\

        RQ5 &
        Machine Learning and Data Engineering generally offer higher compensation. &
        Specialisation may provide greater long-term earning potential. \\

        \bottomrule
    \end{tabularx}
\end{table}

The synthesis presented in Table~\ref{tab:market_synthesis} demonstrates that the five research questions complement one another rather than providing
independent observations. The analyses collectively indicate that successful Python freelancers require a combination of widely applicable programming
skills and domain-specific expertise.

The findings also reveal that market demand cannot be characterised by a single technology or programming library. Instead, employers consistently
expect practical combinations of technologies, particularly Python together with SQL, spreadsheet software, and widely adopted data-analysis libraries.
Furthermore, although Data Analysis represents the largest proportion of the analysed market, the pricing analysis suggests that more specialised fields,
such as Machine Learning and Data Engineering, may offer greater long-term earning potential.

These observations provide the basis for identifying the core technical skills required by the market, which are examined in the following section.

\section{Core Skills Required by the Market} 

The previous sections identified the technologies, programming languages, software tools, and libraries most frequently requested in the analysed
Upwork job postings. Rather than presenting additional statistics, this section synthesizes these findings into a practical hierarchy of technical skills.
The objective is to distinguish between universally required competencies and technologies that are valuable only for specific job types.
This hierarchy provides the foundation for the portfolio project design presented in the next phase.

\begin{table}[H]
\centering
\small
\caption{Skill hierarchy derived from the market analysis.}
\label{tab:skill_hierarchy}

\begin{tabularx}{\textwidth}{l X}
\toprule
\textbf{Skill Level} & \textbf{Technologies} \\
\midrule

Essential &
Python, Pandas, NumPy, SQL, Excel \\

Important &
Power BI, Tableau, Git, Docker, Azure, REST API \\

Specialized &
Machine Learning libraries (Scikit-learn, XGBoost, LightGBM, OpenCV), Selenium, Playwright, OCR frameworks, IBM SPSS, SAS, JavaScript, PostgreSQL,
Microsoft SQL Server, Kubernetes, CI/CD, OpenAI APIs, Claude APIs, Gemini APIs and other domain-specific technologies. \\

\bottomrule
\end{tabularx}
\end{table} 

The market analysis suggests that a relatively small set of technologies forms the foundation of most freelance data-related jobs. Python clearly represents
the central programming language, while Pandas and NumPy constitute the most consistently requested Python libraries. SQL and Excel also appear across
nearly every analysed job category, indicating that database querying and spreadsheet-based data manipulation remain essential practical skills.

A second tier of technologies consists of tools that substantially increase employability but are not required for every job. These include
business intelligence platforms such as Power BI and Tableau, collaborative development tools such as Git, deployment technologies such as Docker,
cloud platforms including Azure, and REST APIs for interacting with external services.

Finally, a number of specialised technologies appear only within particular domains. Machine Learning jobs frequently require dedicated frameworks such as
Scikit-learn, XGBoost or LightGBM, while automation jobs introduce tools such as Selenium and Playwright. Similarly, cloud infrastructure, OCR libraries,
large language model APIs, and orchestration technologies occur only in subsets of the analysed jobs.

This hierarchy indicates that portfolio development should prioritise mastery of the essential technologies before expanding toward specialised domains.
A strong foundation built on Python, SQL, Pandas, NumPy and Excel is therefore expected to maximise applicability across the largest proportion of freelance
opportunities.

\section{Comparison of Python Career Paths} 

The previous sections identified the most common job categories, technologies, required skills, and compensation patterns within the analysed Upwork dataset.
This section combines these findings to provide a qualitative comparison of the four primary Python-related career paths represented in the collected
job postings. Rather than focusing on individual statistics, the objective is to characterise each market segment from the perspective of a
prospective freelance developer.

\begin{table}[H]
\centering
\small
\caption{Comparison of the major Python-related freelance career paths.}
\label{tab:career_comparison}

\begin{tabularx}{\textwidth}{l c c c X}
\toprule
\textbf{Category} &
\textbf{Demand} &
\textbf{Complexity} &
\textbf{Compensation} &
\textbf{Typical Technologies} \\
\midrule

Data Analysis &
High &
Medium &
Medium &
Python, Pandas, NumPy, SQL, Excel, Power BI \\

Data Engineering &
Medium &
High &
High &
Python, SQL, Docker, Azure, REST APIs \\

Machine Learning &
Medium &
Very High &
High &
Python, Scikit-learn, XGBoost, LightGBM, OpenCV \\

Python Automation &
Low &
Medium &
Medium &
Python, Selenium, Playwright, JavaScript \\

\bottomrule
\end{tabularx}

\end{table}

The market analysis reveals that the analysed Python freelance market is composed of several distinct career paths with different characteristics.

Data Analysis represents the largest and most accessible segment of the market. It contains the highest number of job postings and relies primarily on a
relatively stable collection of widely adopted tools, including Python, Pandas, NumPy, SQL, and Excel. Although compensation is generally moderate compared with
more specialised domains, the high demand makes this category particularly attractive for beginners seeking to build practical experience.

Data Engineering occupies a smaller but more technically demanding segment. jobs frequently involve database management, cloud platforms, APIs, and
deployment technologies. While the required technical complexity is noticeably higher, the analysed jobs also indicate a tendency toward larger budgets and
higher hourly rates.

Machine Learning forms the most specialised category within the analysed dataset. In addition to the common Python ecosystem, these jobs require dedicated
machine learning libraries and domain-specific knowledge. Although the number of available jobs is lower than for Data Analysis, the observed budgets suggest
that advanced expertise is generally rewarded with higher compensation.

Python Automation appears as the smallest of the four primary categories. Most jobs focus on browser automation, scripting, or workflow automation using
tools such as Selenium and Playwright. While the technical barrier to entry is lower than for Machine Learning or Data Engineering, the overall number of
opportunities identified in the analysed sample is also considerably smaller.

Overall, the comparison suggests that the different career paths represent different trade-offs between market size, technical complexity, and earning potential.
Data Analysis offers the broadest entry point into the freelance market, whereas Data Engineering and Machine Learning provide opportunities for higher long-term
compensation at the cost of greater technical specialisation. Python Automation occupies a niche position, serving clients with specific automation requirements
rather than general data-related tasks.

\begin{table}[H]
\centering
\small
\caption{Key market insights derived from the analysis.}
\label{tab:key_market_insights}

\begin{tabularx}{\textwidth}{l X}
\toprule
\textbf{Market Insight} & \textbf{Practical Meaning} \\
\midrule

Python is the universal programming language &
Every portfolio project should be implemented primarily in Python. \\

SQL is a core complementary skill &
Database querying is expected across nearly every data-related role. \\

Python alone is rarely sufficient &
Clients typically require combinations of programming, databases and business tools. \\

Specialisation increases earning potential &
Data Engineering and Machine Learning projects generally offer higher compensation. \\

Technology integration is highly valued &
Successful projects combine multiple complementary technologies rather than demonstrating isolated skills. \\

\bottomrule
\end{tabularx}
\end{table}

The findings summarised in Table \ref{tab:key_market_insights} represent the principal conclusions of the market analysis. Rather than describing individual
technologies or job categories, these insights identify the characteristics that consistently distinguish successful freelance projects. Consequently,
these observations provide the design principles for the portfolio roadmap developed in the following phase.

\section{Implications for Portfolio Development}

The analyses presented in the previous sections not only describe the current Python freelance market but also provide clear guidance for
portfolio design. Each freelance category is characterised by a distinct combination of technologies, technical complexity and client expectations.
Consequently, portfolio projects should be designed to reflect the requirements of their target market rather than demonstrating isolated
programming techniques.

\medskip

\textbf{Data Analysis}

A competitive Data Analysis portfolio should demonstrate the ability to analyse real-world datasets and communicate meaningful insights.

\begin{itemize}
    \item use Pyhton as the primary programming language;
    \item demonstrate practical use of Pandas and NumPy for data manipulation;
    \item integrate SQL for data extraction and querying;
    \item include data cleaning and preprocessing;
    \item perfom exploratory data analysis;
    \item produce informative visualisations using libraries such as Matplotlib or Seaborn;
    \item present interpretable conclisons and recommendations based on the analysis;
    \item optionally integrate business intelligence tools such as Power BI or Tableau for interactive dashboards.
\end{itemize}

\medskip

\textbf{Data Engineering}

A competitive Data Engineering portfolio should demonstrate the ability to build reliable and scalable data-processing workflows.

\begin{itemize}
    \item integrate Python with SQL database;
    \item demonstrate automated ETL or ELT piplene;
    \item include data validation and transformation;
    \item interact with external APIs;
    \item demonstrate version control using Git;
    \item include containerisation with Docker where appropriate;
    \item show reprosucible and moudlar project organisation.
\end{itemize}

\medskip
\textbf{Machine Learning}

A competitive Machine Learning portfolio should demonstrate the complete lifecycle of a predictive modelling project.

\begin{itemize}
    \item perfom data preprocessing and feature engineering;
    \item implement machine learning models using Scikit-learn or similar libraries;
    \item compare miltiple modelling approaches;
    \item evaluate model performance using appropriate metrics;
    \item present clear interpretations of the model results;
    \item demonstrate reproducible experimentation and well-structured code.
\end{itemize}

\medskip
\textbf{Python Automation}

A competitive Python Automation portfolio should demonstrate the ability to automate realistic business processes and repetitive tasks.

\begin{itemize}
    \item automate data collection or processing workflows using Python scripts;
    \item interact with web pages or external APIs;
    \item demonstrate browser automation where appropriate;
    \item integrate spreadsheet processing where relevant;
    \item emphasise robutness through logging and error handling;
    \item produce reusable and well-documneted automation scripts.
\end{itemize}

\section[Phase 3 Summary]{Phase Summary}

Phase 3 transformed the quantitative results obtained during the exploratory market analysis into a coherent set of practical market insights.
Rather than examining individual statistics in isolation, the analyses were synthesised to identify the core technical skills required by the market,
compare the major Python-related freelance career paths, and establish category-specific requirements for competitive portfolio projects.

Although the four freelance categories differ considerably in their technical focus, several common characteristics emerged from the market analysis.
Successful portfolio projects should demonstrate practical problem solving, clean software engineering practices, realistic datasets, and the integration
of multiple complementary technologies. At the same time, each career path requires its own combination of domain-specific tools and workflows, which
should be reflected in the portfolio.

Consequently, the market insights developed throughout this phase establish a set of design requirements rather than specific project ideas. These
requirements provide the foundation for the next phase, in which a portfolio roadmap will be constructed by selecting projects that satisfy the
identified market needs while covering the most relevant technologies and skills.